% Source: https://artofproblemsolving.com/community/c3h2057393_vietas_proof
% Stable identity: aops-topic-2057393-q01
Since $p$ is a root of $f(x)=x^2+ax+b$, the factor theorem gives $f(x)=(x-p)g(x)$ where $g$ is a polynomial; comparing degrees and leading coefficients, $g(x)=x-r$ for some number $r$. The roots of $f$ are then exactly $p$ and $r$, so $r=q$ (for a double root $r=p=q$). Thus, for every $x$,
\[x^2+ax+b=(x-p)(x-q)=x^2-(p+q)x+pq .\]
Subtracting, the expression
\[h(x)=\bigl(a+(p+q)\bigr)x+\bigl(b-pq\bigr)\]
equals $0$ for every value of $x$. Putting $x=0$ gives $b-pq=0$. Putting $x=1$ then gives $a+(p+q)=0$. Hence
\[a=-(p+q),\qquad b=pq .\]
This is what "comparing coefficients" means: the coefficient of $x$ on the left is $a$ and on the right is $-(p+q)$; the constant terms are $b$ and $pq$. So the sum of the roots is the opposite of the $x$-coefficient and the product of the roots is the constant term. $\blacksquare$
